% SPDX-FileCopyrightText: 2026 Will Cook
% SPDX-License-Identifier: CC-BY-4.0
%
% Standalone Erdős Problem Note for #1041.  Context is deliberately repeated
% from the sibling notes so that this file remains independently readable.
%
% The public March manuscript is attributable only to the erdosproblems.com
% handle `shtuka`; the PDF contains no real-name author metadata.  The
% bibliography therefore cites that handle and public URL conservatively.
%
% Build from the public paper directory with `make`.
\documentclass[11pt]{article}
% SPDX-FileCopyrightText: 2026 Will Cook
% SPDX-License-Identifier: CC-BY-4.0
%
% Shared preamble for the Erdős Problem Notes: the short, standalone,
% problem-owned manuscripts covering the expansion library `ErdosProblems`.
%
% One note owns one Erdős problem.  Each note repeats whatever context its own
% reader needs, so that a reader who arrives at a single problem never has to
% assemble it from the gateway paper.  Deliberate repetition across notes is the
% design, not an oversight.
%
% Source links resolve against one pinned formal-source commit, declared once
% below.  `scripts/check_problem_note_sources.py` verifies that every authored
% (file, line, declaration) triple names that exact declaration in the pinned
% snapshot, so a link cannot silently rot as later work moves lines.

\usepackage{paper-house-style}
\usepackage{array}
\usepackage{booktabs}
\usepackage{enumitem}
\setlist{topsep=0.35em,itemsep=0.2em}
\usepackage[colorlinks=true,linkcolor=linkinternal,urlcolor=linkexternal,
  citecolor=linkinternal]{hyperref}
\hypersetup{bookmarksnumbered=true,bookmarksopen=true,bookmarksopenlevel=1,
  pdfauthor={Will Cook},pdflang={en-GB}}

% ---------------------------------------------------------------------------
% Pinned formal source.  \commit names the checkpoint every link resolves
% against; the checker reads that snapshot rather than the working tree, so the
% printed coordinates stay correct however later waves move declarations.
% ---------------------------------------------------------------------------
\newcommand{\commit}{e5fd26a6d1b1a346430205eab5385b4c9565d004}
\newcommand{\commitshort}{e5fd26a6d1b1}
\newcommand{\repobase}{https://github.com/wcook04/plectis-lean-erdos249-257/blob/\commit}
\newcommand{\sourceurl}{https://github.com/wcook04/plectis-lean-erdos249-257/tree/\commit}
\newcommand{\PX}{ErdosProblems}

% Declaration names stay inside link targets.  The printed label is either a
% spaced, lower-case reading of the declaration name (\lref) or a short authored
% phrase (\lword); neither prints a module path or a raw Lean identifier.
\ExplSyntaxOn
\cs_new_protected:Npn \noteleanlabel #1
  {
    \tl_set:Nx \l_tmpa_tl { \tl_to_str:n {#1} }
    \regex_replace_all:nnN { _+ } { \c{space} } \l_tmpa_tl
    \regex_replace_all:nnN
      { ([a-z0-9])([A-Z]) }
      { \1\c{space}\2 }
      \l_tmpa_tl
    \tl_set:Nx \l_tmpa_tl { \tl_lower_case:n { \tl_use:N \l_tmpa_tl } }
    \tl_use:N \l_tmpa_tl
  }
\ExplSyntaxOff
\newcommand{\leanlabel}[1]{\noteleanlabel{#1}}
\newcommand{\lref}[3]{\href{\repobase/\PX/#1\#L#2}{\leanlabel{#3}}}
\newcommand{\lword}[4]{\href{\repobase/\PX/#1\#L#2}{#4}}
\newcommand{\lloc}[2]{\href{\repobase/\PX/#1\#L#2}{Lean source}}

% Links into a sibling library, named repository-relative.  The expansion
% library is implicit in \lref/\lword, because that is where problem-owned
% work lands.  Several problems have their headline theorem in the older
% reviewed #249/#257 corpus, and a note that could not name that library
% would have to paraphrase a checked statement instead of linking it.
\newcommand{\mref}[3]{\href{\repobase/#1\#L#2}{\leanlabel{#3}}}
\newcommand{\mword}[4]{\href{\repobase/#1\#L#2}{#4}}
\newcommand{\mloc}[2]{\href{\repobase/#1\#L#2}{Lean source}}
\emergencystretch=2em

\newtheorem{theorem}{Theorem}[section]
\newtheorem{proposition}[theorem]{Proposition}
\newtheorem{lemma}[theorem]{Lemma}
\newtheorem{corollary}[theorem]{Corollary}
\theoremstyle{definition}
\newtheorem{definition}[theorem]{Definition}
\newtheorem{example}[theorem]{Example}
\newtheorem{problem}[theorem]{Problem}
\theoremstyle{remark}
\newtheorem*{remark}{Remark}

\newcommand{\N}{\mathbb{N}}
\newcommand{\Npos}{\mathbb{N}_{>0}}
\newcommand{\Z}{\mathbb{Z}}
\newcommand{\Q}{\mathbb{Q}}
\newcommand{\R}{\mathbb{R}}
\newcommand{\lcm}{\operatorname{lcm}}
\newcommand{\ph}{\varphi}

% The series line printed under every note's title block.
\newcommand{\seriesline}{Erd\H{o}s Problem Notes}

% Production provenance.  The wording is fixed by the corpus provenance
% contract and reproduced verbatim in every manuscript in this repository, so
% the papers cannot drift into different accounts of how they were made.
\newcommand{\noteprovenance}{\provenancenote{\emph{Authorship and AI use.} The
prose, formal proofs, and software in this version were drafted and revised by
large-language-model agents under Will Cook's direction.  Cook set the
objectives and acceptance criteria, reviewed and authorised the manuscript, and
is responsible for its contents.  The agents are production tools, not authors.
See \emph{Statements and declarations}.}}

% The status paragraph every note carries, in its own words, before any result.
% Kernel checking establishes the exact Lean proposition; it does not establish
% that the proposition is the interesting one, that it is new, or that the
% Erdős problem is closed.
\newcommand{\statusboundary}{%
  \emph{Status.}  The problem treated here is open, and this note does not
  close it.  Every statement below marked as checked is a proposition that the
  pinned Lean kernel accepts from the sources this note links to, with no
  \texttt{sorry}, no added axiom, and no unchecked evaluation.  That is a claim
  about the formal statement, not about its mathematical interest, its
  novelty, or the original problem.  The unresolved obligations are named
  exactly, in their own section, and none of the finite computations,
  reductions, or no-go results here removes one of them.}

\renewcommand{\commit}{76b5b0a7ed5da7ebf3b9ed3bfd2fb480b6c38ee0}
\renewcommand{\commitshort}{76b5b0a7ed5d}
\hypersetup{
  pdftitle={Critical-value separation, admissible hubs and Newton-flow geometry for Erdős Problem 1041},
  pdfsubject={Erdős Problem Notes: Problem 1041},
  pdfkeywords={lemniscate, Newton flow, saddle connection, Morse theory, Reeb decomposition, Lean 4},
}

% The prose keeps compact declaration labels; the source appendix supplies the
% commit-pinned links checked by the public release gate.
\newcommand{\pdecl}[1]{\emph{\leanlabel{#1}}}
\newcommand{\dbar}{\mathbb{D}}

\runninghead{Erd\H{o}s \#1041: critical-value separation and admissible hubs}
\title{Critical-Value Separation and Admissible Hubs}
\subtitle{An all-degree sufficient regime, checked dynamical inputs, and the
remaining metric problem for Erd\H{o}s~\#1041}
\author{Will Cook}
\date{August 2026}

\begin{document}
\maketitle
\noteprovenance

\begin{abstract}
\noindent
Let $f(z)=\prod_i(z-z_i)$ be monic with roots in the closed unit disc.  Erd\H{o}s,
Herzog, and Piranian ask whether two roots can always be joined by a curve of
length less than $2$ inside $\{|f|<1\}$.  The global assertion remains open.

The strongest all-degree sufficient regime proved here is controlled by
critical-value separation.  Let $c$ be a simple critical point with
$v=f(c)\ne0$.  If every other critical point $d$ satisfies
$|1-f(d)/v|\ge S>1$, then the square-root-resolved inverse branch through $c$
gives a root connector of length at most
\[
 2|v|^{1/n}(1+S)^{1/n}
       \sqrt{\log\!\frac{S}{S-1}}
\]
inside $\{|f|\le|v|\}$.  Hence the connector is shorter than
$2|v|^{1/n}$ whenever
$(1+S)^{2/n}\log(S/(S-1))<1$.  Exact uniform choices are $S=4$ for
$n\ge3$, $S=3$ for $n\ge4$, and $S=2$ for $n\ge7$.  For roots in the open
unit disc these hypotheses prove Erd\H{o}s~\#1041 for the polynomial.  The
analytic proof remains ordinary mathematics; Lean separately checks degree
monotonicity, the three exact thresholds, and the squared-length consumer.

A second all-degree result is sharp rather than merely sufficient: for
collinear zeros of diameter $D$, one adjacent-root segment lies below
$C_n(D/2)^n$, where
$C_n=(2^{n-1}\cos^n(\pi/(2n)))^{-1}$, and equality occurs for affine
Chebyshev configurations.  We also close the primitive sparse quintic family
and every translated cubic quotient-fibre family of degree $3q$, $q\ge2$.
For each family the exact Lean-checked selector is displayed separately from
the ordinary normalization, fibre, moment, and path assembly.

This result exposes rather than hides the surviving frontier: it gives no
control when critical values form a near tie, and it excludes multiple
saddles.  Those are precisely the clustered regimes still requiring a grouped
contour, an admissible-hub selector, or another argument.  Complementary exact
inputs sharpen that boundary.  At every critical point, two nearest-root
occurrences have total Euclidean distance at most $2$; Lean checks the
reciprocal critical balance and disk inverse-square estimate used by the
ordinary polynomial-level argument, but these estimates alone do not prove
containment.  An exact quintic has a nearest straight spoke that leaves
$\{|f|<1\}$, and an exact cubic has an escaping midpoint on every straight
root-pair segment.

Along the complex Newton field $-f/f'$, the value satisfies
$f(z(t))=e^{-t}f(z(0))$. Lean checks this value equation, the resulting
positive-ray obstruction to finite saddle connections, and the finite
translation and root-retention interfaces.

Outside the separated regime, the remaining metric problem is one of
admissible-hub selection.  On the ray-separated dense class, one must find a
critical hub $c$ whose two inverse-ray arms stay in the relevant component and
satisfy
\[
  \min_{c\ \mathrm{admissible}} L(c)\le2.
\]
At degree five, a critical value at most $1/M_5$, where
\[
 M_5=(1-t_*)(1+t_*)^3\sqrt{16t_*^2-4t_*+1},
 \qquad t_*={5\over16}+{3\sqrt{105}\over80},
 \qquad 1/M_5=0.2760461\ldots,
\]
would give two contained nearest-root spokes of total length at most $2$.
The near-Fekete reduction removes radial deficits from this containment
question, but neither the general hub selection nor its degree-five form is
proved. The threshold and no-go witnesses are ordinary proofs, exact
certificates, or computations, not additional kernel-checked theorems.
Coefficient genericity, component stability, planar topology, and metric
gluing remain open.
\end{abstract}

\begin{paperresult}{What the analysis isolates}
\textbf{Exact inputs.} Critical-point balance gives two nearest-root
distances with total at most $2$, while the Newton flow transports $f$
exactly by $f(z(t))=e^{-t}f(z(0))$.  Exact examples show why straight spokes
and root-pair segments do not solve the containment problem.  \textbf{Sharp
degree-five target.} A critical value at most $1/M_5$ would suffice.
\textbf{Open boundary.} The missing theorem is an admissible critical hub
whose inverse-ray arms stay in the right component and meet the metric bound.
\end{paperresult}

\externalreviewobject{The paper proves an all-degree critical-value-separation
regime by a resolved inverse branch, the area formula, P\'olya's
area--capacity inequality, and an exact coefficient estimate.  It proves a
sharp all-degree collinear theorem, a complete primitive sparse quintic
theorem, and translated cubic quotient-fibre theorems in every degree $3q$.
It also checks exponential Newton-value decay, positive-ray collision
interfaces, finite translation avoidance, and quantified root retention, and
records the current near-Fekete residual.}{The separation and three
solved-family assemblies are ordinary mathematics, not kernel-checked
authority.
It excludes near-tied critical values and multiple saddles.  The checked
dynamical and perturbative inputs do not repair the global topology and metric
gluing in those complementary strata, so Problem~\#1041 remains open.  The
degree-five target and the no-go witnesses are boundaries around the open
theorem, not a hidden closure.}{Comparator
selects the finite-family small-translation theorem and the quantified
root-retention theorem, three solved-family kernels, and the numerical kernel
of critical-value separation.  Each formal endpoint routes to the exact paper
result and boundary it supports; the covering-space and area argument and the
frontier section remain ordinary evidence classes.}

The two selected perturbation interfaces have the following exact boundaries.
For a finite type $\iota$, an injective family of critical values
$c\colon\iota\to\mathbb C$, and every $\varepsilon>0$, the row
\texttt{exists\_small\_translation\_separating\_arguments} supplies a shift
of norm less than $\varepsilon$ such that every translated value is nonzero
and every two distinct translated values lie on different positive rays.
Separately, for a monic split polynomial of positive degree whose roots have
norm at most $\rho$, the row
\texttt{constant\_perturbation\_roots\_in\_unitDisk} requires $\rho\ge0$,
$\varepsilon>0$, the margin
\[
 ((\operatorname{natDegree}f+1)\varepsilon)^{1/\operatorname{natDegree}f}
   +\rho<1,
\]
and a shift of norm less than $\varepsilon$; it then places every root of
$f+\operatorname{C}(\text{shift})$ in the open unit disc.  These are finite
perturbation and stability interfaces only: neither row supplies the missing
global topology or the length-$2$ gluing argument.

\section{The problem}\label{sec:problem}

\begin{problem}[Erd\H{o}s \#1041]\label{res:problem}
Let $f(z)=\prod_{i=1}^{n}(z-z_i)$ be monic with $z_i\in\dbar$, the open unit
disc.  Show that two of the roots can be joined by a curve of length less than
$2$ lying in the open lemniscate $E=\{z\in\mathbb{C}:|f(z)|<1\}$.
\end{problem}

\noindent Numbering and current status follow Bloom's Erd\H{o}s problem
catalogue \cite{bloom}.
The problem is open.  The original source is Problem~5 on printed p.~139 of
Erd\H{o}s--Herzog--Piranian \cite[p.~139]{ehp1958}; the preceding paragraph records the
known input that one component of the lemniscate contains at least two zeros.

Recent work on polynomial lemniscates separates component counts from metric
path questions. Ghosh and Ramachandran characterize the number of components
through critical points and critical values \cite[Lemma~7]{ghosh2023}; for the
binomial family $z^n-a$, the condition $|a|<1$ puts the filled unit lemniscate
in the connected regime. Connectedness alone gives no path-length bound.
Our dynamical terminology is also standard: Sutherland calls
$\dot z=-f(z)/f'(z)$ the continuous Newton flow and observes that $f(z(t))$
moves on a straight radial line \cite[p.~42]{sutherland1992}. We use this
value-space identity, not a claim that the global trajectory graph is a tree.

Two recent manuscripts are relevant.  The 48-page manuscript posted by
\texttt{shtuka} on 24 March 2026
\cite[Theorem~1, p.~1]{march2026} claims the unrestricted
statement.  Its Proposition~12 (p.~16, with proof continuing through p.~30)
supplies the spanning-tree decomposition used in the final proof.  The defect
was located publicly in the problem's discussion thread: on 25 March 2026 Tao
observed that the invocation of Lemma~8 there is unjustified and that the flow
lines need not organise into connected trees, and on 26 March 2026 the
manuscript's author agreed that the statement of Proposition~12 itself, not
only its printed proof, is incorrect, and set the strategy aside.  Section~\ref{sec:gap}
records an independent diagnosis of the same failure through the local
three-ended saddle model, together with possible repairs.  No counterexample
to the proposition is exhibited there.  Pendyala's independent June 2026
preprint \cite[Thm.~1, p.~1]{june2026} proves the degree-four case.  The quartic theorem does
not close the problem.

We study the Newton flow whose trajectories foliate the lemniscate.

\statusboundary

\begin{table}[t]
\centering
\footnotesize
\begin{tabular}{@{}
  >{\raggedright\arraybackslash}p{0.34\linewidth}
  >{\raggedright\arraybackslash}p{0.20\linewidth}
  >{\raggedright\arraybackslash}p{0.38\linewidth}@{}}
\toprule
\papertablehead
Statement & Status & Exact boundary \\
\midrule
Erd\H{o}s \#1041 & Open & No proof is claimed. \\
Critical-value separation & Ordinary all-degree theorem; Lean-checked numerical kernel & A simple nonzero hub and $(1+S)^{2/n}\log(S/(S-1))<1$; covering and area inputs stay ordinary; near ties and multiple saddles are excluded. \\
Newton value equation $w'=-w$ & Checked & Away from critical points, along any trajectory tangent to $-f/f'$. \\
Exponential first integral & Checked & $\tfrac{d}{dt}\bigl(e^{t}f(z(t))\bigr)=0$. \\
Ray separation of critical values & Checked (consumer form) & Endpoints of a finite connection share one oriented ray; distinct rays exclude a connection. \\
Ray-collision locus & Checked & $\beta=(ra-b)/(1-r)$, $r>0$, $r\ne1$: one real parameter per pair. \\
Quartic case & Cited & Proved in \cite[Thm.~1, p.~1]{june2026}; does not extend to general degree. \\
Translated quartic quotient fibres & Ordinary theorem with a Lean-checked metric kernel & For $f(z)=P((z-h)^q)$, $P$ monic quartic and $q\ge2$; Pendyala supplies the quartic geometry, while Lean checks the root-lift density, exact primitive, and strict endpoint budget. \\
Signed-moment cyclic tetranomials & Lean-checked two-index safe-spoke theorem & For an indexed finite root family of $g(w)=w^m+aw^r+bw^s+c$, an exact signed $L^2$ moment budget selects two distinct indices whose complete spokes are safe; distinct root values require injectivity of the indexing map. \\
Concyclic zeros with $2\rho^n\le1$ & Ordinary proof; finite exact and numerical checks & Distinct-root theorem; not Lean checked; the unrestricted concyclic case remains open. \\
Unrestricted proof of \cite[Theorem~1, p.~1]{march2026} & Proof gap & Proposition~12 uses a false three-ended local saddle block; located publicly by Tao (25 March 2026), conceded by the author at statement level (26 March 2026).  No counterexample is exhibited. \\
Constant-translation ray separation and root retention & Checked & After critical-value injectivity, arbitrary small ray avoidance and an explicit unit-disc margin. \\
Coefficient perturbation and slack stability & Open & Must first create injective critical values and preserve the component, collars and length budget. \\
Reeb decomposition and length fan-in & Open & The two surviving producers. \\
Random search to degree $10$ & Verified finite instances & Upper bounds for sampled configurations only. \\
\bottomrule
\end{tabular}
\end{table}

\section{An all-degree critical-value separation theorem}
\label{sec:critical-value-separation}

The first result is the strongest direct parent-theorem regime in this note.
It is not a statement about a sampled family or a fixed degree: the condition
is an exact inequality on the other critical values after normalising one
simple first-merge hub.

\begin{theorem}[critical-value separation]
\label{res:critical-value-separation}
Let $P$ be a polynomial of degree $n\ge3$ whose leading coefficient has
modulus one.  Suppose
\[
 P(0)=1,\qquad P'(0)=0,\qquad P''(0)\ne0,
\]
and, for some $S>1$, every other critical point $d\ne0$ satisfies
\[
 |1-P(d)|\ge S.
\]
If $Z$ is either local solution at the saddle of
\[
 P(Z(\xi))=1-\xi^2,\qquad Z(0)=0,
\]
then $Z$ continues holomorphically and injectively to
$|\xi|<\sqrt S$.  Its endpoints $Z(-1)$ and $Z(1)$ are distinct roots, and
their resolved inverse-ray connector satisfies
\[
 \int_{-1}^{1}|Z'(\xi)|\,d\xi
 \le 2(1+S)^{1/n}\sqrt{\log\!\frac{S}{S-1}}.       \tag{4}
\]
In particular the connector has length strictly below $2$ whenever
\[
 (1+S)^{2/n}\log\!\frac{S}{S-1}<1.                 \tag{5}
\]
\end{theorem}

\begin{proof}
The simple saddle has local form $P(z)=1+Az^2+O(z^3)$ with $A\ne0$, so the
substitution $P(Z(\xi))=1-\xi^2$ resolves it into two local holomorphic
sheets.  Normalize the algebraic curve
$X=\{(\xi,z):P(z)=1-\xi^2\}$.  A finite ramification point of its projection
to the $\xi$-plane can occur only at a critical point $d$ of $P$, where
$\xi^2=1-P(d)$.  The node over $(0,0)$ is replaced in the normalization by
the two resolved local points, and the separation hypothesis puts every
other ramification point on or outside $|\xi|=\sqrt S$.  Polynomial
properness makes the normalized projection finite and proper.  Its inverse
image over $D_{\sqrt S}$ is therefore an unramified finite cover.  Since the
disc is simply connected, every connected cover component maps
biholomorphically to it; in particular either resolved local sheet continues
as a single holomorphic branch $Z$ throughout the disc.

This branch is also injective as a map into the $z$-plane.  Indeed,
$Z(\xi_1)=Z(\xi_2)$ implies $\xi_1^2=\xi_2^2$.  The involution
$\iota(\xi,z)=(-\xi,z)$ exchanges the two distinct normalized points over
$(0,0)$ and hence exchanges their two cover components.  If
$\xi_2=-\xi_1\ne0$ and $Z(\xi_1)=Z(\xi_2)$, one component would meet its
$\iota$-image; connected cover components are disjoint, so they would have to
coincide, contradicting the two distinct points over $0$.  Thus
$\xi_1=\xi_2$.

Write $Z(\xi)=\sum_{k\ge1}a_k\xi^k$.  For
$1<R<\sqrt S$, injectivity and the area formula give
\[
 \pi\sum_{k\ge1}k|a_k|^2R^{2k}
   =\operatorname{Area}(Z(D_R)).                    \tag{6}
\]
The identity $P(Z(\xi))=1-\xi^2$ places the image in
$\{|P|<1+R^2\}$.  Since the leading coefficient of $P$ has modulus one,
P\'olya's area--capacity inequality \cite[printed pp.~280--282]{polya1928}
yields
\[
 \operatorname{Area}(Z(D_R))
 \le \operatorname{Area}\{|P|<1+R^2\}
 \le \pi(1+R^2)^{2/n}.                              \tag{7}
\]
Cauchy--Schwarz applied to (6)--(7) now gives
\[
 \sum_{k\ge1}|a_k|
 \le (1+R^2)^{1/n}
      \sqrt{\sum_{k\ge1}\frac1{kR^{2k}}}
 = (1+R^2)^{1/n}\sqrt{\log\!\frac{R^2}{R^2-1}}.   \tag{8}
\]
Termwise integration on $[-1,1]$ gives
$\int_{-1}^{1}|Z'|\le2\sum_{k\ge1}|a_k|$.  Letting
$R\nearrow\sqrt S$ proves (4).  Injectivity makes $Z(-1)$ and $Z(1)$
distinct, while $P(Z(\xi))=1-\xi^2\in[0,1]$ on the real segment proves the
claimed containment.  Squaring (4) gives (5).
\end{proof}

\begin{corollary}[exact convenient thresholds]
\label{res:critical-value-thresholds}
Condition~\emph{(5)} holds for
\[
 (n,S)=(n,4)\quad(n\ge3),\qquad
 (n,3)\quad(n\ge4),\qquad
 (n,2)\quad(n\ge7).
\]
If $f$ is monic with roots in the open unit disc, $c$ is a simple critical
point with $v=f(c)\ne0$, and every other critical point satisfies the
corresponding inequality $|1-f(d)/v|\ge S$, then two roots of $f$ are joined
inside $\{|f|<1\}$ by a curve of length strictly below $2$.
\end{corollary}

\begin{proof}
For $S=4$ and $n\ge3$,
\[
 5^{2/n}\log(4/3)
 \le5^{2/3}\log(1+1/3)<5^{2/3}/3<1,
\]
where the last radical inequality is $25<27$.  For $S=3$ and $n\ge4$,
$4^{2/n}\log(3/2)\le2\log(1+1/2)<1$.  For $S=2$ and $n\ge7$, the positive
exponential series gives
$\exp(7/10)>12013/6000>2$, hence $\log2<7/10$, while
$9\cdot7^7<10^7$ is exactly $(3^{2/7}\cdot7/10)^7<1$.

For the polynomial $f$, set $z=c+\rho e^{i\theta}w$ with
$\rho=|v|^{1/n}$ and choose $\theta$ so that $P(w)=f(z)/v$ has
unit-modulus leading coefficient.  The theorem gives a connector of length
strictly below $2|v|^{1/n}$ inside $\{|f|\le|v|\}$.  For each displayed
threshold $S\ge2$, separation gives $|f(d)|\ge(S-1)|v|\ge|v|$ at every other
critical point, so $|v|=\mu$ is the least critical-value modulus.  If the
roots lie in a disc of radius $R<1$, the resultant identity and Fekete's
Vandermonde bound give, with critical points counted with multiplicity,
\[
 \mu^{n-1}
 \le \prod_{f'(d)=0}|f(d)|
 =\frac{|\operatorname{Disc}(f)|}{n^n}
 =\frac{\prod_{i<j}|z_i-z_j|^2}{n^n}
 \le R^{n(n-1)}.
\]
Thus $|v|=\mu\le R^n<1$, so both the length and containment are strict at
the target scale.
\end{proof}

\paragraph{Evidence and exact boundary.}
\label{bdry:critical-value-separation}
The analytic continuation, monodromy, area formula and P\'olya inequality in
Theorem~\ref{res:critical-value-separation} are ordinary mathematics.  The
companion symbolic replay checks the branch-value algebra, all three exact
constant chains, and the nonempty boundary example
$P(z)=1-3z^2+z^3$, whose other critical point is $2$ and satisfies
$1-P(2)=4$.  The companion Lean module
	exttt{FirstMergeCriticalValueSeparation} checks the degree monotonicity,
all three exact coefficient thresholds, and the abstract implication from
the squared connector estimate to length below two.  It does not formalize
the analytic continuation, univalence, area formula, or P\'olya inequality,
and therefore does not turn the complete analytic proof into a Lean theorem.

The boundary is structural, not cosmetic.  This method says nothing when a
second critical value enters the resolved disc, and it assumes the selected
saddle is simple.  It therefore leaves the near-tie and multiple-saddle
strata, makes no sharpness claim for the convenient thresholds, and does not
prove an unrestricted admissible-hub selector, a COVER theorem, or
Erd\H{o}s~\#1041.  No novelty or priority claim is made here for the use of
P\'olya's global area inequality; the contribution asserted is the displayed
composed sufficient regime and its exact constants.  P\'olya is cited only
for that global input, not for the square-resolved covering and coefficient
assembly.

\section{Three exact solved polynomial families}
\label{sec:solved-polynomial-families}

The next three theorems are unconditional subcases of the path problem, not
proxies for the unrestricted conjecture.  They are ordered by mathematical
signal rather than by the date of their formalization: first an all-degree
sharp theorem with equality configurations, then a complete degree-five
sparse family, and finally an infinite tower of translated quotient fibres.
In each case a displayed finite proposition is the exact Comparator-selected
Lean endpoint.  The ensuing normalization and path theorem is ordinary
mathematics and is stated separately.

\subsection{Sharp collinear zeros}
\label{subsec:sharp-collinear-chebyshev}

Put
\[
 r_n=\cos\frac{\pi}{2n},
 \qquad C_n=\frac{1}{2^{n-1}r_n^n}.
\]

\begin{theorem}[checked Chebyshev endpoint]
\label{prop:sharp-collinear-chebyshev-comparator}
Let $m\ge0$, let $p\in\mathbb R[X]$ be monic of degree $m+2$, and let
\[
 -1<c_0<\cdots<c_m<1,\qquad |c_i|\le1.
\]
Suppose $p(-1)=p(1)=0$ and
$p(c_i)p(c_{i+1})<0$ for $0\le i<m$.  Then
\[
 \min_{0\le i\le m}|p(c_i)|\le C_{m+2}.
\]
\end{theorem}

This is exactly the statement selected as
\pdecl{Erdos249257.ExternalVerification1041SolvedFamilies.SharpCollinear.existsPeakLeComparisonBound}.
It is the kernel-facing alternation endpoint, not yet the geometric theorem.

\begin{theorem}[sharp collinear diameter theorem]
\label{thm:sharp-collinear-diameter}
Let $f$ be a monic polynomial of degree $n\ge2$ whose zero occurrences are
collinear, and let $D$ be their diameter.  Some two adjacent zero occurrences
are joined by a segment of length at most $D$ on which
\[
 |f(z)|\le
 \frac{(D/2)^n}{2^{n-1}\cos^n(\pi/(2n))}.          \tag{9}
\]
The constant in \emph{(9)} is best possible in every degree.  Equality is
attained by affine images of the zeros of $T_n$ whose extreme zeros have
distance $D$.
\end{theorem}

\begin{proof}
A repeated zero gives the constant path, so assume the zero values are
distinct.  A rigid motion sends their line to the real axis, their midpoint
to the origin, and their extremes to $\pm D/2$.  With $R=D/2$, the normalized
polynomial
\[
 q(w)=R^{-n}e^{-in\theta}f(m+Re^{i\theta}w)
\]
is monic with real zeros
$-1=y_1<\cdots<y_n=1$, and
$|f(m+Re^{i\theta}w)|=R^n|q(w)|$.

Compare $q$ with the monic endpoint-normalized Chebyshev polynomial
\[
 q_*(x)=\frac{T_n(r_nx)}{2^{n-1}r_n^n}.
\]
Both $q$ and $q_*$ vanish at $\pm1$, and $|q_*|\le C_n$ on
$[-1,1]$.  For each gap $[y_i,y_{i+1}]$, choose $c_i$ at which
$|q|$ is maximal.  The signs of $q(c_i)$ alternate.  If every gap maximum
were larger than $C_n$, then $q-q_*$ would have the same alternating signs
as $q$ at the $c_i$.  It would therefore have a zero between each consecutive
pair $c_i,c_{i+1}$, as well as the two zeros $\pm1$.  These are $n$ distinct
zeros, whereas the leading terms cancel and
$\deg(q-q_*)\le n-1$; moreover $q-q_*$ is nonzero at every $c_i$.
This contradiction selects a gap on which $|q|\le C_n$.  Scaling back proves
\emph{(9)}.

For sharpness, take
\[
 y_k=\frac{\cos((2k-1)\pi/(2n))}{r_n},
 \qquad 1\le k\le n.
\]
These are the zeros of $q_*$, have extremes $\pm1$, and every adjacent gap
contains a scaled Chebyshev extremum where $|q_*|=C_n$.  No smaller universal
constant can work.
\end{proof}

\begin{corollary}[collinear Erd\H{o}s case]
\label{cor:collinear-erdos-1041}
If the zero occurrences of a monic polynomial lie on one line in the open
unit disc, two of them are joined by a curve of length strictly below $2$
inside $\{|f|<1\}$.
\end{corollary}

\begin{proof}
Their diameter satisfies $D<2$.  Also $C_n\le1$ for $n\ge2$, with equality
only at $n=2$, so \emph{(9)} is strictly below one.  The selected segment has
length at most $D<2$.
\end{proof}

\paragraph{Formal boundary.}
\label{rem:sharp-collinear-formal-boundary}
Lean checks sign preservation, the alternating root-count mechanism, the
scaled Chebyshev endpoint and uniform bound, and
Theorem~\ref{prop:sharp-collinear-chebyshev-comparator}.  It does not check
the rigid normalization, the existence of the gap maxima, transport back to
the root line, or the equality-node locations used in
Theorem~\ref{thm:sharp-collinear-diameter}.  These are ordinary steps.  The
comparison with classical Chebyshev alternation is explicit, and no
literature-priority claim is made for the fixed-diameter adjacent-gap
assembly.

\subsection{The primitive sparse quintic}
\label{subsec:primitive-sparse-quintic}

The next selector is finite but unusually rigid.  For $0<r<2$ and
$0\le s_i\le1$, $x_i^2\le s_i$, put
\[
 E_i=s_i^4(s_i+r^2+2rx_i).
\]

\begin{theorem}[checked two-tail energy selector]
\label{prop:primitive-quintic-two-tail-energy-selector}
Suppose
\[
 \sum_{i=0}^4x_i=-r,\qquad
 \sum_{i=0}^4(2x_i^2-s_i)=r^2,\qquad
 \sum_{i=0}^4(4x_i^3-3s_ix_i)=-r^3.              \tag{10}
\]
Then at least one of the ten pairs $0\le i<j\le4$ satisfies
$E_i<1$ and $E_j<1$.
\end{theorem}

This is the exact finite conclusion selected as
\pdecl{Erdos249257.ExternalVerification1041SolvedFamilies.primitiveQuintic_twoStrictTailEnergies}.
It selects two distinct indices.  It does not assert that the corresponding
complex root values are distinct.

\begin{theorem}[primitive sparse quintic]
\label{thm:primitive-quintic-two-tail}
Let
\[
 p(z)=z^5+az^4+bz+c
\]
and suppose its five zero occurrences $w_0,\ldots,w_4$ lie in the closed unit
disc.  At least two distinct indices satisfy
\[
 |bw_i+c|\le1.                                    \tag{11}
\]
If $a\ne0$, two indices can be chosen with strict inequalities.  If $a=0$,
every index satisfies \emph{(11)}, and equality holds exactly when
$|w_i|=1$.

For open-disc zeros, two zero occurrences are joined inside $\{|p|<1\}$ by a
curve of length below $2$: use the two radial spokes through $0$ when their
values are distinct, and the constant path when the selected occurrences
have the same value.
\end{theorem}

\begin{proof}
Rotate so that $a=re^{i\phi}$ becomes $r\ge0$ and write
$z_i=e^{-i\phi}w_i$.  The missing $z^3$ and $z^2$ coefficients and Newton's
identities give
\[
 \sum z_i=-r,\qquad \sum z_i^2=r^2,\qquad
 \sum z_i^3=-r^3.                                 \tag{12}
\]
At a zero,
\[
 |bw_i+c|=|w_i|^4|w_i+a|.
\]
Thus, for $x_i=\Re z_i$ and $s_i=|z_i|^2$, the squared tail is exactly
$E_i$ and \emph{(12)} becomes \emph{(10)}.

For completeness, the selector is driven by the harmonic extension
\[
 H_r(x,s)=(-r/2-x)(1-x)^2
 +(1-s)\left(1-\frac r4-\frac{3x}{4}\right).
\]
The three moments give the exact sum
\[
 \sum_{i=0}^4H_r(x_i,s_i)=5-2r,                  \tag{13}
\]
while $H_r\le4-2r$ throughout the unit disc.  If $E_r(x,s)\ge1$, then
\[
 5(1-s)\le2r(x+r/2).
\]
After putting $d=x+r/2$ and $u=1-x$, the remaining estimate is the exact
sum-of-squares identity
\[
 \frac1{31}-P
 =\frac25\left(d-\frac{38}{31}
               +\frac14(u-\frac3{31})\right)^2
  +\frac{31}{40}\left(u-\frac3{31}\right)^2\ge0,
\]
which yields $H_r(x,s)\le2/31$ for every unsafe tail.  If four tails were
unsafe, \emph{(13)} would give
\[
 5-2r\le4-2r+\frac8{31}<5-2r,
\]
a contradiction.  This proves the two-index conclusion.  The case $r=0$ is
the direct identity $|bw_i+c|=|w_i|^5$.

Finally, at a zero $w$,
\[
 p(tw)=(1-t)c+(t-t^4)(bw+c)-t^4(1-t)w^5.
\]
For $0\le t\le1$ the three nonnegative coefficients sum to $1-t^5$.
The tail bound, $|w|\le1$, and $|c|\le1$ therefore give the radial
unit-sublevel path.  Under the open-disc hypothesis the relevant bounds and
the total spoke length $|w_i|+|w_j|$ are strict.
\end{proof}

\paragraph{Formal boundary.}
\label{rem:primitive-quintic-formal-boundary}
Lean checks the five-index moment hypotheses, the harmonic cap, the
sum-of-squares unsafe estimate, and the ten-pair conclusion in
Theorem~\ref{prop:primitive-quintic-two-tail-energy-selector}.  The rotation,
derivation of Newton moments from the roots, root-tail identity, Abel
decomposition, path assembly, and any passage from distinct indices to
distinct complex values remain ordinary mathematics.  No priority claim is
made for the resulting complete sparse family.

\subsection{Translated cubic quotient fibres}
\label{subsec:translated-cubic-quotient-fibres}

\begin{theorem}[checked cubic safe spoke]
\label{lem:cubic-safe-root-spoke}
If $r,s,v\in\mathbb C$ have modulus below one, at least one
$u\in\{r,s,v\}$ satisfies
\[
 \left|(tu-r)(tu-s)(tu-v)\right|\le1
 \qquad(0\le t\le1).
\]
\end{theorem}

This is exactly
\pdecl{Erdos249257.ExternalVerification1041SolvedFamilies.cubic_safeRootSpoke}.

\begin{theorem}[translated cubic quotient fibres]
\label{thm:translated-cubic-quotient-fibres}
Let $q\ge2$, $h\in\mathbb C$, $P$ be monic cubic, and
\[
 f(z)=P((z-h)^q).
\]
If every zero of $f$ lies in the open unit disc and $f$ has at least two
distinct zero values, then two zeros are joined through $h$ by a two-segment
path of length below $2$ inside $\{|f|<1\}$.  Equivalently this closes the
coefficient family
\[
 (z-h)^{3q}+A(z-h)^{2q}+B(z-h)^q+C
\]
in every degree $3q\ge6$.
\end{theorem}

\begin{proof}
Write $P(w)=(w-r)(w-s)(w-v)$ and assign to $r$ the charge
$A_r=\Re(r\overline{s+v})$, cyclically.  The exact sum is
\[
 A_r+A_s+A_v=|r+s+v|^2-(|r|^2+|s|^2+|v|^2)>-3,
\]
so one charge, say $A_r$, exceeds $-1$.  For $0\le t\le1$, AM--GM and the
distance-square identity give
\[
 |tr-s|\,|tr-v|<1+t+t^2.
\]
Consequently
\[
 |P(tr)|<(1-t)(1+t+t^2)=1-t^3<1                 \tag{14}
\]
away from the root endpoint.

Fix a $q$th root $y$ of a quotient root and a primitive $q$th root of unity
$\zeta$.  The full fibre consists of $h+y\zeta^k$.  Since every fibre point
lies in the open unit disc,
\[
 \frac1q\sum_{k=0}^{q-1}|h+y\zeta^k|^2
   =|h|^2+|y|^2<1.                                \tag{15}
\]
Thus $|y|<1$.  For the safe quotient root from \emph{(14)}, two distinct
fibre points satisfy
\[
 f(h+ty\zeta^k)=P(t^qr),
\]
and their two spokes through $h$ have total length $2|y|<2$.  If the selected
quotient root is zero, choose a nonzero quotient root and use the direct
estimate $|P(ts)|<2t(1-t)\le1/2$; if none exists, $f$ has only one distinct
zero value, contrary to the hypothesis.
\end{proof}

\paragraph{Formal boundary.}
\label{rem:cubic-quotient-formal-boundary}
Lean checks the charge identity and selection, the two-distance envelope,
the cubic product bound, the safe-spoke disjunction in
Theorem~\ref{lem:cubic-safe-root-spoke}, and an exact quartic falsifier for
this charge method.  It does not check the finite $q$-fibre construction,
root-of-unity average \emph{(15)}, zero-root fallback, pullback, or geometric
path assembly.  These are ordinary steps, and no literature-priority claim is
made.

\paragraph{Shared assurance boundary.}
\label{bdry:solved-polynomial-families}
The three Comparator declarations route respectively to
Theorems~\ref{prop:sharp-collinear-chebyshev-comparator},
\ref{prop:primitive-quintic-two-tail-energy-selector}, and
\ref{lem:cubic-safe-root-spoke}.  They support, but are strictly weaker than,
the ordinary assembled Theorems~\ref{thm:sharp-collinear-diameter},
\ref{thm:primitive-quintic-two-tail}, and
\ref{thm:translated-cubic-quotient-fibres}.  None of these solved scopes
proves the unrestricted problem, and none licenses a novelty, Comparator,
NanoDa, Palomar, or publication claim beyond its separately recorded receipt.

\section{The current frontier: what survives near Fekete configurations}
\label{sec:frontier}

The August~29 public research update changes the useful first reading of this
problem.  It does not add a proof of Erd\H{o}s~\#1041; it identifies the
geometric regime in which the remaining difficulty is concentrated and kills
several attractive but false shortcuts.  The source of the update is pinned to
commit
\href{https://github.com/wcook04/plectis-lean-erdos249-257/tree/f214a6b45528dc5eefe20ffadc35f2e981627d4c/research_corpus/Erdos1041}{\texttt{f214a6b4}}
and the dated synthesis is
\href{https://github.com/wcook04/plectis-lean-erdos249-257/blob/f214a6b45528dc5eefe20ffadc35f2e981627d4c/research_corpus/Erdos1041/FRONTIER.md}{\texttt{FRONTIER.md}}.
The distinctions below are part of the result: a theorem, a refutation, a
certificate, and a measurement do not have interchangeable force.

\subsection*{Near-Fekete stability removes the radial distraction}

Write $a_i=\rho_i u_i$, with $|u_i|=1$, and put
$D=|\operatorname{disc}(f)|/n^n$.  The quantitative
Fekete--Hadamard estimate in
\href{https://github.com/wcook04/plectis-lean-erdos249-257/blob/f214a6b45528dc5eefe20ffadc35f2e981627d4c/research_corpus/Erdos1041/NearFeketeRadialAngularSplit.md}{\texttt{NearFeketeRadialAngularSplit.md}}
says that, when $D\ge1-\eta$ and
$\eta\le1/(80n^2)$,
\[
  1-\rho_i^2\le \frac{n\eta}{n-1},\qquad
  |a_i-a_j|\ge \frac{2-2\sqrt\eta-n\eta}{n-1},
\]
and the roots lie within $7\sqrt\eta$ of a rotated regular $n$-gon, after a
bijection.  The mechanism is a normalised Vandermonde Gram determinant:
near equality in Hadamard's inequality forces both the radii and the angular
separation to be near extremal.  This is an ordinary mathematical result in
the public research corpus, not a new Lean declaration.

The associated radial monotonicity theorem is the more important reader
conclusion.  Under the stronger near-Fekete condition
$\eta\le1/(10n^4)$ used by this corollary, replacing the radii by one can only
increase the values along the corresponding radial spoke.  Thus
radial deficits help containment; the unresolved near-Fekete problem is
angular.  On exact regular-gon directions this becomes unconditional for
$n\le6$: arbitrary radii in the stated band give
\[
 |f(s\,\omega\zeta^i)|\le1-s^n
 \quad(0\le s\le\rho_i),
\]
so every two-radii path is contained.  The good spoke in the general
on-circle averaging identity may move with $s$, however; that is why an
existence statement at one level does not supply a fixed pair of roots.

\subsection*{A bounded-radius concyclic class}

There is nevertheless a clean theorem for a different on-circle mechanism.
Let $f$ be monic of degree $n\ge3$ with distinct zeros on a circle of radius
$\rho$.  If $2\rho^n\le1$, two adjacent zeros are joined by their straight
chord, whose length is at most $2\rho\sin(\pi/n)<2$, and the chord lies in
$\{|f|<1\}$.  (If a zero is repeated, the short-connection conclusion is
immediate; the distinct case is the substantive one.)

After normalisation to
the unit circle, self-inversive realification and alternation against $z^n-c$
select a zero-free gap on which the polynomial modulus is at most $2$; a
harmonic normal-derivative argument transfers this arc bound strictly to the
chord.  The complete ordinary proof is
\href{https://github.com/wcook04/plectis-lean-erdos249-257/blob/f214a6b45528dc5eefe20ffadc35f2e981627d4c/research_corpus/Erdos1041/ConcyclicAlternation.md}{\texttt{ConcyclicAlternation.md}}.

The argument is an ordinary proof outside Lean.  The
\href{https://github.com/wcook04/plectis-lean-erdos249-257/blob/f214a6b45528dc5eefe20ffadc35f2e981627d4c/research_corpus/Erdos1041/scripts/check_erdos1041_concyclic_exact_witness.py}{exact-rational checker}
checks finitely many load-bearing identities and configurations, while the
\href{https://github.com/wcook04/plectis-lean-erdos249-257/blob/f214a6b45528dc5eefe20ffadc35f2e981627d4c/research_corpus/Erdos1041/scripts/check_erdos1041_concyclic_alternation.py}{numerical checker}
is regression and stress-test evidence.  The arc constant $2$ is sharp on the
regular $n$-gon, so this chord argument does not reach radii tending to $1$;
the unrestricted concyclic case and Erd\H{o}s~\#1041 remain open.

\subsection*{Cyclic trinomial fibres}

There is also a positive structured family in every cyclic lift degree.  Fix
integers $1\le r\le m$ and $q\ge1$, and consider the translated polynomial
\[
 f(z)=(z-h)^{qm}+a(z-h)^{qr}+c.
\]
Writing $w=(z-h)^q$ reduces the root equation to
$w^m+aw^r+c=0$.  At such a quotient root the middle coefficient can be
eliminated exactly: for $0\le u\le1$,
\[
 u^mw^m+au^rw^r+c
   =(1-u^r)c-(u^r-u^m)w^m.
\]
The two coefficients on the right are nonnegative and have total at most one.
Lean therefore proves that, if $|w|<1$ and $|c|<1$, the complete radial spoke
is strictly contained in the open unit lemniscate.  This is an all-root
statement; no Vieta-small root selection is needed for the containment bound.

For a nontrivial cyclic fibre, two selected displacements $y_1,y_2$ with
$|y_1|,|y_2|<1$ also satisfy $|y_1|+|y_2|<2$.  The formal source checks the
factorization, strict spoke estimate, and this metric budget.  It does not yet
check the ordinary finite argument selecting a suitable quotient root and two
distinct members of its fibre, lifting both quotient spokes, or assembling
their union as a path in the relevant component.  Thus the Lean theorem is the
load-bearing analytic kernel for this translated cyclic-trinomial family, not
the complete path theorem and not a proof of unrestricted Erd\H{o}s~\#1041.

\subsection*{Coefficient-controlled cyclic tetranomial spokes}

The first additional lower monomial admits an exact Abel-tail criterion.  Let
\[
  g(w)=w^m+aw^r+bw^s+c,\qquad m>r>s\ge1,
\]
assume all roots of $g$ lie in the open unit disk, and let $w_1,w_2$ be roots
of the two smallest moduli.  If
\[
  |c|+|b|\,|w_2|^s<1,
\]
then the complete radial spokes from $w_1$ and $w_2$ to the origin lie in
$\{|g|<1\}$, and their broken line has length strictly below $2$.  The simpler
coefficient condition
\[
  |b|+|c|\le1
\]
makes every root spoke safe; the coefficient $a$ is unrestricted.

The mechanism is visible in one identity.  At a root $w$ and for
$0\le u\le1$,
\[
\begin{split}
 g(uw)={}&(1-u^s)c
  -(u^s-u^r)(aw^r+w^m)
  -(u^r-u^m)w^m,\\
 aw^r+w^m={}&-(c+bw^s).
\end{split}
\]
The three scalar coefficients are nonnegative and sum to $1-u^m$.  Thus the
root-dependent budget $|c|+|b||w|^s<1$, together with $|w|<1$ and $|c|<1$,
puts every vector in the Abel decomposition strictly inside the unit ball.
Lean checks this factorization and the resulting strict spoke theorem under
the exact weak exponent hypotheses $1\le s\le r\le m$; it also checks the
coefficient-only corollary above.

There is a stronger, genuinely two-index formal theorem.  Let $S$ index a
finite family of roots, $N=|S|\ge2$, and
\[
  M=\sum_{i\in S}w_i^s.
\]
Lean proves the exact signed energy identity
\[
 \sum_{i\in S}|c+bw_i^s|^2
 =N|c|^2+|b|^2\sum_{i\in S}|w_i^s|^2
   +2\mathop{\rm Re}(\overline c\,bM).
\]
Consequently, if every $|w_i|<1$, $|c|<1$, and
\[
 N\bigl(|b|^2+|c|^2\bigr)
   +2\mathop{\rm Re}(\overline c\,bM)<N-1,
\]
then two distinct indices $i,j\in S$ satisfy
$|c+bw_i^s|<1$ and $|c+bw_j^s|<1$.  Feeding those inequalities into the
Abel decomposition proves, in Lean, that both complete root spokes lie
strictly in $\{|g|<1\}$.  The signed cross term is essential: this interface
can certify configurations outside the coefficient-only triangle
$|b|+|c|\le1$.  The formal hypotheses do not require $i\mapsto w_i$ to be
injective, so distinct indices need not denote distinct root values; the
ordinary two-root consequence requires a root enumeration without repetition.

For $q\ge2$ and $f(z)=g((z-h)^q)$, the ordinary regular-fibre mean-square
identity puts every quotient root inside the unit disk.  Lifting two selected
quotient spokes gives a path through $h$ whose two displacement lengths sum
to less than $2$.  The new Lean theorem assumes the finite root family, its
signed moment $M$, and the displayed budget; it does not derive that budget
from arbitrary tetranomial coefficients, prove that the supplied family is a
complete root multiset, perform cyclic-fibre lifting, or construct the final
path object.  Without the signed-moment, root-dependent, or coefficient-only
budget, the tetranomial case remains open; this family does not solve
unrestricted Erd\H{o}s~\#1041.

\subsection*{Translated quartic quotient fibres}

Pendyala's degree-four theorem can also be lifted through every nontrivial
cyclic power.  Let $P$ be a monic quartic, let $q\ge2$, and set
\[
  f(z)=P((z-h)^q).
\]
If all listed zeros of $f$ lie in the open unit disk, averaging the squared
moduli over each complete $q$-point fibre shows that every quotient root
$w$ of $P$ satisfies $|w|<1$.  Pendyala's chord-or-radial proof then supplies
the quotient geometry.  The extra issue is metric: a short chord in the
$w$-plane need not lift isometrically through $y\mapsto y^q$.

Put $\alpha=1/q$.  Along a quotient chord whose supporting line has distance
$d$ from the origin, the root-lift density is controlled pointwise by
\[
  \bigl(\sqrt{d^2+x^2}\bigr)^{\alpha-1}
    \le x^{\alpha-1}\qquad(x>0),
\]
and Lean checks the exact primitive
\[
  \alpha\int_0^A x^{\alpha-1}\,dx=A^\alpha.
\]
Splitting at the perpendicular foot, and at the origin when the chord crosses
it, yields the ordinary power-map chord-lift estimate
\[
  \operatorname{length}(\widetilde{[a,b]})
    \le |a|^{1/q}+|b|^{1/q}.
\]
The formal source also checks the strict endpoint consumer: if
$0\le a,b<1$, $\alpha>0$, and a candidate length $L$ is at most
$a^\alpha+b^\alpha$, then $L<2$.  Thus both Pendyala branches survive the
power-map lift, giving the asserted short path for every translated quartic
quotient fibre, in each degree $4q\ge8$.

The attribution and proof boundary are exact.  Pendyala proves the quartic
geometric theorem and its four-point radial lemma.  The local Lean module
checks the antitone density inequality, its integral, the strict powered
endpoint budget, and the final length fan-in.  It does not formalize
Pendyala's geometric lemma, the continuous covering-space construction of the
root lift, or the ordinary chord/radial case assembly.  This is a structured
all-scale family, not a proof of unrestricted Erd\H{o}s~\#1041.

\subsection*{The degree-five target is now explicit}

The degree-five analysis sharpens the two-nearest-spoke threshold.  If a
critical point $c$ satisfies $|f(c)|\le1/M_n$, then both straight segments to
the two nearest roots remain in $\{|f|\le1\}$ and their total length is at
most $2$.  The envelope is obtained by maximising the exact product along the
second spoke; at degree five,
\[
 M_5=(1-t_*)(1+t_*)^3\sqrt{16t_*^2-4t_*+1},
 \qquad
 t_*={5\over16}+{3\sqrt{105}\over80},
 \qquad
 {1\over M_5}=0.2760461\ldots.
\]
This is an 11-fold improvement over the earlier deep-low threshold, but it
does not assert that every quintic has a critical value below this level.

The exact remaining degree-five statement is therefore not ``improve the
constant'' but the following selection problem:
\begin{quote}
For every monic quintic with roots in the closed unit disc, there is a
critical point $c$ with $|f(c)|\le1$ and two roots $a,b$ such that
$|f|\le1$ on $[c,a]\cup[c,b]$ and
$|c-a|+|c-b|\le2$.
\end{quote}
The public file calls this (SPOKE-5).  Together with the existing
two-segment reduction it would settle degree five, but (SPOKE-5) is not
proved.  Coverage experiments put the residual in rapid, nearly simultaneous
merging near the regular pentagon: the generic sampled families were covered,
whereas the surviving band has every named merge threshold active and
$D$ near $1$.  Those percentages locate the work; they are not a theorem.

At the same degree, the terminal connected component admits a sharper area
identity.  If $K_t=\{|f|\le t\}$ is connected and the centred exterior map is
\[
 \psi(\zeta)=t^{1/n}\zeta+a_0+
 \sum_{k\ge1}a_k\zeta^{-k},
\]
then Gr\"onwall's identity gives
\[
 \operatorname{Area}(K_t)=\pi\left(t^{2/n}-
 \sum_{k\ge1}k|a_k|^2\right),
 \qquad
 a_1=-\frac{c_{n-2}}{n\,t^{1/n}}
 \quad\text{after centring}.
\]
It replaces a P\'olya upper bound at that terminal node, but extending it to
proper components and connecting it to the parent theorem remains open.  The
degree-five statements and their evidence classes are recorded in
\href{https://github.com/wcook04/plectis-lean-erdos249-257/blob/f214a6b45528dc5eefe20ffadc35f2e981627d4c/research_corpus/Erdos1041/Degree5AssemblyAndSharpenedCuts.md}{\texttt{Degree5AssemblyAndSharpenedCuts.md}}.

\subsection*{The no-go boundaries force hub selection}

The same frontier is useful because it closes off three misleading readings.
First, the strict-argmin first-merge hub need not have two-arm length below
$2$.  A degree-five tie-locus witness, moved into the open unit disc, has
\[
 L(c_*)=2.0573432753\ldots>2,
\]
with a 50-digit two-instrument reconstruction and a margin about 97 times
that of the earlier degree-four witness.  This is computational certificate
evidence, not an exact rational theorem; the parent remains viable because a
different hub at the same configuration supplies the relevant pair.

Second, the aggregate estimate
$\sum_cL(c)\le2(n-1)R$ is refuted at degree four on an open violating region
and at one degree-five witness.  Its algebraic factor
\[
 \sum_{k=1}^{n-1}|f(c_k)|^{1/n}\le(n-1)R
\]
survives as a separate conjecture, but the metric factor is not automatic.
These two boundaries, including the distinction between the surviving
algebraic half and the dead aggregate, are in
\href{https://github.com/wcook04/plectis-lean-erdos249-257/blob/f214a6b45528dc5eefe20ffadc35f2e981627d4c/research_corpus/Erdos1041/MinimalHubArmBudgetRefutation.md}{\texttt{MinimalHubArmBudgetRefutation.md}}
and
\href{https://github.com/wcook04/plectis-lean-erdos249-257/blob/f214a6b45528dc5eefe20ffadc35f2e981627d4c/research_corpus/Erdos1041/SeparatrixAggregateReduction.md}{\texttt{SeparatrixAggregateReduction.md}}.

Third, the origin is not a uniform near-Fekete hub.  An exact rational
quintic with roots scaled by $999999/1000000$ has four of five origin spokes
escaping $\{|f|\le1\}$ at explicit rational sample points, even though
$1-D=3.19991\cdot10^{-4}$.  Consequently no fixed $\eta$-neighbourhood of
the Fekete locus can force two contained origin spokes in degree five.  The
mechanism is angular: the first several Fourier modes can point against
different spokes.  The certificate is exact for the displayed witness; the
claim that the phenomenon persists throughout a limiting family is
computational.  The full witness and replay route are in
\href{https://github.com/wcook04/plectis-lean-erdos249-257/blob/f214a6b45528dc5eefe20ffadc35f2e981627d4c/research_corpus/Erdos1041/NearFeketeRadialAngularSplit.md}{\texttt{NearFeketeRadialAngularSplit.md}}.

\subsection*{The surviving carrier and the actual open endpoint}

These refutations leave a selector as the surviving carrier.  On the
ray-separated dense class, the public frontier records the canonical open row
\[
  \min_{c\ \mathrm{admissible}} L(c)\le2,
\]
where admissibility includes the two inverse-ray arms being contained in the
relevant lemniscate component.  Its measured values remain below $2$ at the
refuting configurations, and the existing attachment and lower
semicontinuity reductions would turn this row into the parent theorem.  No
proof of the row is currently available.  In particular, the nonzero hub
that rescues the origin-spoke witnesses is evidence for the selector, not a
universal construction.

The honest current endpoint is therefore a containment selector with three
separate obligations: choose the hub, keep both arms inside the component,
and control the metric fan-in.  The Lean module cited elsewhere in this note
checks the Newton value equation, ray-collision avoidance, and root-retention
inputs; the new frontier files above are ordinary proof, exact-certificate,
and computational research records, not formal authority for (SPOKE-5), the
no-go generalisations, or the parent theorem.  Keeping those authority
boundaries visible is part of making the frontier reusable.

\section{The Newton value equation}\label{sec:newton}

Away from the critical set, define the complex Newton field
\[
  N(z)=-\frac{f(z)}{f'(z)} .
\]
Let $z(t)$ be differentiable with $z'(t)=N(z(t))$, and put $w(t)=f(z(t))$.

\begin{theorem}[value equation]\label{res:value}
$w'(t)=-w(t)$.
\end{theorem}

\noindent The computation is one line: $w'=f'(z)\,z'=f'(z)\cdot(-f(z)/f'(z))=-f(z)=-w$.
The kernel checks it as \pdecl{newtonFlow_value_hasDerivAt}, together with the
differential form of the first integral,
\pdecl{newtonFlow_scaledValue_hasDerivAt_zero}:
\[
  \frac{d}{dt}\Bigl(e^{t}f(z(t))\Bigr)=0,
  \qquad\text{equivalently}\qquad
  f(z(t))=e^{-t}f(z(0)) .
\]
Observe what this says about the geometry.  The value moves radially inward at
exponential rate and never changes argument.  The lemniscate $\{|f|<1\}$ is
therefore forward-invariant, and the flow lines are exactly the preimages of
rays from the origin.

\begin{corollary}[ray separation]\label{res:ray}
The endpoints of any finite Newton-flow connection lie on the same oriented ray
from $0$.  Consequently, if two critical values lie on distinct positive rays,
no Newton-flow saddle connection joins the corresponding critical points.
\end{corollary}

\noindent This is checked in consumer form: the kernel accepts the implication
from the hypothesis of distinct rays to the absence of a connection.  It is the
statement the topology needs, and it is a genuine sharpening of the criterion
used in the literature.

\section{Arguments, not moduli}\label{sec:arguments}

It is tempting to arrange a generic perturbation so that the critical values
are pairwise distinct, or that their moduli are pairwise distinct, and to
conclude that saddle connections are excluded.  Neither is enough.

Two distinct critical values can lie on one ray, and two critical values with
distinct moduli certainly can: the ray records the argument, and the modulus is
exactly the coordinate the flow contracts.  By Corollary~\ref{res:ray} the
invariant that excludes connections is the argument.  What a perturbation must
therefore achieve is pairwise distinct critical-value \emph{arguments}, which is
a condition on $n-1$ points modulo the circle rather than on their positions in
the plane.

The cost of that condition is also checked, and it is small.

\begin{theorem}[ray-collision locus]\label{res:locus}
Let $a\ne b$ be complex.  Every common translation $\beta$ for which $a+\beta$
and $b+\beta$ lie on the same positive ray has the form
\[
  \beta=\frac{ra-b}{1-r},
  \qquad r\in\mathbb{R}_{>0},\ r\ne1 .
\]
\end{theorem}

\noindent So each pair of critical values contributes a one-real-parameter
forbidden locus in the translation plane, given in closed form
(\pdecl{translated_samePositiveRay_parameterization}).  A finite union of such
loci has empty interior, which is the shape one wants for an avoidance
argument.  Turning that into a perturbation of $f$ is not immediate: the
translation model must be replaced by an actual perturbation of the roots that
keeps them inside $\dbar$ and preserves the length slack.  That is the first
open producer of \S\ref{sec:open}.

\section{A proof gap in the unrestricted argument}\label{sec:gap}

The March manuscript's Proposition~12 claims the following load-bearing
statement.  For $u=-\log|f|$, a connected component $V$ of $\{u>c\}$ carrying
$m\ge2$ simple zeros, and the stated regularity and Morse hypotheses, it
constructs, for every $\varepsilon>0$, an embedded spanning tree
$G_\varepsilon\subset V$ with
\[
 \operatorname{len}(G_\varepsilon)
 \le\frac1{2\pi}\int_{2\alpha}^{\infty}P_V(t)\,dt+\varepsilon.
\tag{5.1}\label{eq:prop12-bound}
\]
The final theorem uses this estimate, so the issue below cannot be bypassed by
calling the proposition auxiliary.

At an interior index-one critical point $p$, the proof invokes a Morse chart
\[
 u=\mu+x^2-y^2
\]
and replaces the saddle by a three-ended neighbourhood having one connected
lower cross-section and two connected upper cross-sections.  That local model
is false as written.  Because $p\in V$ and $V$ is open, a sufficiently small
closed disc around $p$ lies entirely in $V$.  In that disc the full Morse chart
has four sectors: two components of $u>\mu$ and two components of $u<\mu$.
A global component argument cannot delete one local sector from a disc already
contained in $V$.

This diagnoses a proof step, not the proposition's statement.  A repair might
cut an adjoining regular annulus along a separatrix or regular flow arc before
forming the block, retain a four-pronged saddle neighbourhood and change the
assembly, or replace the local construction by the ray-cut decomposition
proposed below.  Any repair must prove that its connector cost can be made
arbitrarily small uniformly in the attachment points and that the repaired
blocks still assemble to an embedded tree satisfying~\eqref{eq:prop12-bound}.
The shorter descriptions of the same three-ended block do not repair the
four-sector topology.

Corollary~\ref{res:ray} supplies one independent input for a different route:
distinct critical-value arguments exclude saddle-to-saddle Newton
connections.  It does not itself prove the compact planar decomposition,
classify all orbit endpoints or provide the metric gluing estimate.  Those are
separate problems below.

\section{Finite evidence}\label{sec:finite}

A search was run over random monic polynomials with roots in the unit disc.  For
each sample the region $\{|f|<1\}$ was rasterised and shortest grid paths were
computed between every pair of roots.  The best upper bounds obtained were
\[
  \begin{array}{lrr}
    \text{degree } 5, & 500 \text{ trials}: & 1.1052648928\\
    \text{degree } 6, & 1500 \text{ trials}: & 0.8450414343\\
    \text{degree } 8, & 1500 \text{ trials}: & 0.6203916714\\
    \text{degree } 10, & 1500 \text{ trials}: & 0.4303640486 .
  \end{array}
\]
No counterexample candidate was found, and the measured values sit well below
the threshold $2$.  These are grid distances for the sampled configurations:
upper bounds on those samples, not bounds over the family.  The apparent
decrease with degree is a property of the sample, and we draw no conjecture from
it.
Future searches should target named failure modes---near-degenerate saddles,
thin necks, boundary-critical configurations and almost-connected
separatrices---and report those diagnostics.  Raster paths remain candidate
finders, never continuous certificates.

\section{Complements and further questions}\label{sec:open}

The dependency chain has five separate gates.  A proof, a minimally corrected
hypothesis set, or an explicit polynomial or planar counterexample is a
decisive answer to any one of them.

\subsection*{1. Repair or refute the saddle block}

\begin{problem}[corrected local saddle assembly]\label{prob:saddle1041}
Under the hypotheses of Proposition~12, replace the invalid three-ended local
model by a four-pronged block or by a block formed after a specified annular
cut.  Prove that for every $\eta>0$ its connector has total Euclidean length
below $\eta$, uniformly over the selected attachment points, and that the
corrected blocks assemble to an embedded tree satisfying
\[
 \operatorname{len}(G_\varepsilon)
 \le\frac1{2\pi}\int_{2\alpha}^{\infty}P_V(t)\,dt+\varepsilon;
\]
or give a polynomial or harmonic planar counterexample to that statement.
\end{problem}

The full-disc model $x^2-y^2$, an annulus in which lower branches rejoin, two
saddles joined by a separatrix and simultaneous saddle levels are mandatory
tests.  Repeating the one-lower/two-upper assertion does not answer the
problem.

\subsection*{2. The compact ray-cut decomposition}

Let $p$ have simple roots, simple nonzero critical points and let
$u=-\log|p|$.  For regular values $c<T$, take
\[
 M_{c,T}=\overline{V\cap\{c<u<T\}},
\]
and assume explicitly that this is a compact genus-zero surface, its lower
boundary is one smooth Jordan curve, its upper boundary consists of $m$ smooth
root curves, all interior critical points are nondegenerate index-one saddles,
the normalised gradient field
\[
 X=\frac{\nabla u}{|\nabla u|^2}=-\frac{p}{p'}
\]
is transverse to the level boundaries, and no maximal $X$-trajectory has two
saddle endpoints.

\begin{problem}[finite strip decomposition after ray cuts]\label{prob:reeb1041}
For every $\eta>0$, construct disjoint Morse neighbourhoods $N_j$ with
$\sum_j\operatorname{diam}N_j<\eta$ and a finite set of complete separatrix or
regular-flow cuts so that every component of the complement is flow-diffeomorphic
to a rectangle
\[
 [a_S,b_S]\times[0,1],\qquad u(\Phi_S(t,s))=t,
\]
with connected level sections and no uncut annular component.  Give an explicit
finite bound $E(s,m)$ for the number of strips, or exhibit the minimal missing
hypothesis or a counterexample.
\end{problem}

No particular formula such as $2s+1$ is presumed.  Boundary tangencies,
simultaneous levels, branch reunion through an annulus and non-Hausdorff orbit
spaces must be handled rather than suppressed.

\subsection*{3. Metric fan-in without losing the coefficient}

For a strip $S$, write
\[
 \Gamma_t^S=S\cap\{u=t\},\qquad
 P_S(t)=\mathcal H^1(\Gamma_t^S),
\]
and let $k_S$ be its number of root ends.  The flux identity
\[
 \int_{\Gamma_t^S}|\nabla u|\,ds=2\pi k_S
\]
gives the average trajectory estimate
\[
 \int_{\Gamma_{t_0}^S}\operatorname{len}(\gamma_x)\,d\mu_{t_0}(x)
 =
 \frac1{2\pi k_S}\int_{a_S}^{b_S}P_S(t)\,dt.
\]

\begin{problem}[additive-error strip gluing]\label{prob:metric1041}
Assuming Problem~\ref{prob:reeb1041}, select trajectories in all strips and
connect them through the saddle and root neighbourhoods so that, for every
$\eta>0$, the resulting embedded tree contains all $m$ roots and obeys
\[
 \operatorname{len}(G)
 \le\frac1{2\pi}\int_{2\alpha}^{\infty}P_V(t)\,dt+\eta.
\tag{6.1}\label{eq:metric-fanin1041}
\]
The total saddle, annular-cut and root-cap cost must be below $\eta$ without a
multiplicative loss in $1/(2\pi)$.
\end{problem}

For the final strict inequality, use the actual collar slack
\[
 q=\frac1{2\pi}\int_{\alpha}^{2\alpha}P_V(t)\,dt>0
\]
and give budgets that keep the perturbation, tree error and transfer cost below
fixed fractions of $q$.  Independently choosing a shortest trajectory in each
strip is not enough unless the attachment mismatch is controlled.

\subsection*{4. Coefficient perturbation and stability}

The constant-translation stage is no longer open.  Once a finite critical-value
family is injective, Lean proves an arbitrarily small translation making every
value nonzero and pairwise positive-ray separated
(\pdecl{exists_small_translation_separating_arguments}).  It also proves the
explicit root-retention estimate; a shift below $\varepsilon$ keeps all roots
in the unit disc when
\[
 ((n+1)\varepsilon)^{1/n}+\rho<1
\]
(\pdecl{constant_perturbation_roots_in_unitDisk}).  A constant translation
cannot separate initially equal critical values.

\begin{problem}[two-stage generic perturbation with slack]\label{prob:perturb1041}
For fixed root discs, compact collar $K$, regular levels
$\alpha/2,\alpha,2\alpha,5\alpha/2$ and collar slack $q>0$, prove that an
arbitrarily small
\[
 g_{\lambda,\beta}(z)=f(z)+\lambda z+\beta
\]
can be chosen so that:

\begin{enumerate}
\item $f+\lambda z$ has simple relevant critical points and injective complex
critical values;
\item the checked constant translation $\beta$ makes those values nonzero and
pairwise ray-separated;
\item no critical point enters the protected collar or truncation boundary;
\item the relevant component remains in $K$, contains exactly the corresponding
perturbed roots and has slack $q_g>q/2$; and
\item root displacement and straight-line transfer back to the original roots
consume less than a prescribed fraction of $q/m$.
\end{enumerate}

If the one-coefficient perturbation $\lambda z$ cannot ensure all five
properties, give the smallest additional lower-coefficient direction that can,
or an explicit obstruction.
\end{problem}

Finite planar avoidance and the subsequent constant translation must not be
relisted as missing; coefficient genericity, component stability and slack
stability are the open content.

\subsection*{5. The global Newton-flow claim ceiling}

\begin{problem}[relative global Newton-flow theorem]\label{prob:globalflow1041}
On a compact regular band
\[
 M_{a,b}=\{z:a\le-\log|p(z)|\le b\},
\]
prove real-time existence and the identities
\[
 p(z(t))=e^{-(t-t_0)}p(z(t_0)),\qquad
 u(z(t))=u(z(t_0))+t-t_0
\]
throughout every maximal orbit; classify both limiting endpoints among the
regular boundary, root ends and finite saddle set; and exclude recurrent,
accumulating and non-Hausdorff orbit phenomena.  Under pairwise ray separation
of the critical values, decide whether the flow is Morse--Smale relative to the
boundary or whether the strongest valid conclusion is only absence of
saddle-to-saddle connections.
\end{problem}

The checked algebra proves the pointwise value equation and the endpoint-ray
consumer.  It does not supply global solution theory or an orbit-space graph.
A positive stronger theorem must give the graph and a finite edge bound; a
negative answer should exhibit the simplest ray-separated polynomial carrying
the remaining pathology.

Erd\H{o}s~\#1041 remains open.  The source now publicly verifies the Newton
kernel, finite ray avoidance and quantitative constant-translation root
control; the coefficient perturbation, corrected planar decomposition and
metric gluing remain the exact unresolved producers.

\section*{Statements and declarations}

Lean does not check the exposition, citation choices, or interpretation. This
manuscript cites Lean only for the formal statements and proofs that the pinned
kernel accepts. The checked
core is the Newton value equation, the exponential first integral, the consumer
form of ray separation, the finite planar-avoidance theorem, quantitative
constant-translation root retention, and the ray-collision parameterisation.  The
decomposition and length statements of \S\ref{sec:open} are not proved.  The
diagnosis of Proposition~12 in \S\ref{sec:gap} concerns its printed local
saddle construction; it does not refute the proposition's statement.  The
search results of \S\ref{sec:finite} are computations.

\appendix
\section{Guide to the formal sources}\label{app:sources}

The public \texttt{ErdosProblems.Erdos1041.NewtonFlowRaySeparation} module
contains the checked source for this note.  The search of \S\ref{sec:finite} is
\texttt{scripts/search\_counterexample.py} in the source package.  The
declaration table below is pinned to the shared formal-source commit used
throughout this problem-note series.

\begin{itemize}[leftmargin=*]
\item \lref{Erdos1041/NewtonFlowRaySeparation.lean}{34}{newtonFlowVector}
\item \lref{Erdos1041/NewtonFlowRaySeparation.lean}{38}{derivative_mul_newtonFlowVector}
\item \lref{Erdos1041/NewtonFlowRaySeparation.lean}{50}{newtonFlow_value_hasDerivAt}
\item \lref{Erdos1041/NewtonFlowRaySeparation.lean}{64}{newtonFlow_scaledValue_hasDerivAt_zero}
\item \lref{Erdos1041/NewtonFlowRaySeparation.lean}{77}{SamePositiveRay}
\item \lref{Erdos1041/NewtonFlowRaySeparation.lean}{80}{samePositiveRay_refl}
\item \lref{Erdos1041/NewtonFlowRaySeparation.lean}{84}{samePositiveRay_symm}
\item \lref{Erdos1041/NewtonFlowRaySeparation.lean}{92}{samePositiveRay_trans}
\item \lref{Erdos1041/NewtonFlowRaySeparation.lean}{107}{translated_samePositiveRay_parameterization}
\item \lref{Erdos1041/NewtonFlowRaySeparation.lean}{127}{realAffineLine}
\item \lref{Erdos1041/NewtonFlowRaySeparation.lean}{130}{realAffineLine_eq_affineSpan}
\item \lref{Erdos1041/NewtonFlowRaySeparation.lean}{147}{isClosed_realAffineLine}
\item \lref{Erdos1041/NewtonFlowRaySeparation.lean}{152}{dense_compl_realAffineLine}
\item \lref{Erdos1041/NewtonFlowRaySeparation.lean}{162}{exists_small_avoiding_finite_realAffineLines}
\item \lref{Erdos1041/NewtonFlowRaySeparation.lean}{179}{samePositiveRay_imp_mem_realAffineLine}
\item \lref{Erdos1041/NewtonFlowRaySeparation.lean}{197}{exists_small_translation_separating_arguments}
\item \lref{Erdos1041/NewtonFlowRaySeparation.lean}{230}{norm_lt_one_of_near_root}
\item \lref{Erdos1041/NewtonFlowRaySeparation.lean}{257}{constant_perturbation_root_near_original}
\item \lref{Erdos1041/NewtonFlowRaySeparation.lean}{287}{constant_perturbation_roots_in_unitDisk}
\item \lref{Erdos1041/NewtonFlowRaySeparation.lean}{306}{samePositiveRay_of_real_exp_decay}
\item \lref{Erdos1041/NewtonFlowRaySeparation.lean}{315}{no_newtonConnection_of_not_samePositiveRay}
\end{itemize}

\begin{thebibliography}{9}
\bibitem{polya1928} G.~P\'olya,
  \emph{Beitrag zur Verallgemeinerung des Verzerrungssatzes auf mehrfach
  zusammenh\"angende Gebiete}, Sitzungsber. Preuss. Akad. Wiss.,
  Phys.-Math. Kl. (1928), 228--232 and 280--282.
\bibitem{bloom} T.~F. Bloom, \emph{Erd\H{o}s Problems}, problem~1041.
  \url{https://www.erdosproblems.com/1041}
\bibitem{ehp1958} P.~Erd\H{o}s, F.~Herzog, and G.~Piranian,
  \emph{Metric properties of polynomials}, J. Analyse Math. \textbf{6}
  (1958), 125--148. \url{https://doi.org/10.1007/BF02790232}
\bibitem{ghosh2023} S.~Ghosh and K.~Ramachandran,
  \emph{Number of Components of Polynomial Lemniscates: A Problem of Erd\H{o}s,
  Herzog, and Piranian}, arXiv:2312.13673v1 (2023).
  \url{https://arxiv.org/abs/2312.13673}
\bibitem{sutherland1992} S.~Sutherland,
  \emph{Bad Polynomials for Newton's Method}, in B.~Bielefeld and M.~Lyubich
  (eds.), \emph{Conformal Dynamics Problem List}, Stony Brook IMS preprint
  (1992), 42--44.
  \url{https://www.math.stonybrook.edu/preprints/ims92-7.pdf}
\bibitem{march2026} \texttt{shtuka},
  \emph{A Short Path Joining Two Zeros Inside a Polynomial Lemniscate},
  manuscript posted 24 March 2026, 48~pp.
  \url{https://shtuka123.github.io/1041/main.pdf}.
  The file at this URL has since been replaced by a shorter partial version
  that no longer contains Proposition~12; the durable public record of the
  March version, its defect, and the author's 26 March 2026 concession is the
  discussion thread at
  \url{https://www.erdosproblems.com/forum/thread/1041}.
\bibitem{june2026} V.~S. Pendyala,
  \emph{A Degree-Four Lemniscate Path Theorem},
  arXiv:2606.24875v1 (2026).
  \url{https://arxiv.org/abs/2606.24875},
  doi:\href{https://doi.org/10.48550/arXiv.2606.24875}
  {10.48550/arXiv.2606.24875}.
\end{thebibliography}

\end{document}
